
      \documentclass[fleqn]{article}      
      \usepackage{amsmath}
      \usepackage{fontspec}
      \usepackage{kotex} % 한국어 지원  

      \begin{document}      
      쉬운 객관식 질문: \\\\  
1. 함수란 무엇인가요? \\\\
   a) 두 수를 더하는 것 \\\\
   b) 입력에 대해 출력이 정해진 규칙 \\\\
   c) 도형을 그리는 것 \\\\

2. 다음 중 함수의 예는 무엇인가요? \\\\
   a) $y = 2x + 3$ \\\\
   b) $x + y = 5$ \\\\
   c) $x^2 + y^2 = 1$ \\\\

3. $f(x) = x + 1$에서 $f(2)$의 값은? \\\\
   a) 1 \\\\
   b) 2 \\\\
   c) 3 \\\\

중간 객관식 질문: \\\\  
1. 함수의 그래프는 어떤 모양을 가질 수 있나요? \\\\
   a) 직선 \\\\
   b) 곡선 \\\\
   c) 둘 다 \\\\

2. $f(x) = x^2$의 그래프는 어떤 형태인가요? \\\\
   a) 직선 \\\\
   b) 포물선 \\\\
   c) 원 \\\\

3. 다음 중 함수의 정의에 맞지 않는 것은? \\\\
   a) 각 입력에 대해 오직 하나의 출력이 존재함 \\\\
   b) 하나의 입력이 여러 개의 출력을 가질 수 있음 \\\\
   c) 입력과 출력의 관계가 규칙적임 \\\\

어려운 객관식 질문: \\\\  
1. 함수의 역함수란 무엇인가요? \\\\
   a) 입력과 출력을 바꾼 함수 \\\\
   b) 출력이 항상 0인 함수 \\\\
   c) 입력을 제곱한 함수 \\\\

2. $f(x) = 3x + 2$의 역함수는 무엇인가요? \\\\
   a) $f^{-1}(x) = \frac{x-2}{3}$ \\\\
   b) $f^{-1}(x) = \frac{3x - 2}{1}$ \\\\
   c) $f^{-1}(x) = 2x + 3$ \\\\

3. $f(x) = x^3$의 출력이 27일 때, 입력 $x$는 무엇인가요? \\\\
   a) 3 \\\\
   b) 9 \\\\
   c) 27 \\\\

주관식 문제: \\\\  
쉬움: \\\\
1. 함수의 정의를 설명하시오. \\\\
2. $f(x) = x + 5$에서 $f(3)$의 값을 구하시오. \\\\
3. $f(x) = 2x$일 때, $f(4)$의 값을 구하시오. \\\\

중간: \\\\
1. $f(x) = x^2 + 1$의 그래프는 어떤 모양인가요? \\\\
2. 함수의 입력과 출력의 관계를 설명하시오. \\\\
3. $f(x) = 3x - 4$에서 $f(2)$의 값을 구하시오. \\\\

어려움: \\\\
1. 함수 $f(x) = 2x + 1$과 $g(x) = x^2$의 합성함수 $h(x) = f(g(x))$를 구하시오. \\\\
2. $f(x) = 4x - 1$의 역함수를 구하시오. \\\\
3. $f(x) = 5x^2 + 3x + 2$의 최대값을 구하시오. \\\\

      
      \end{document} 
  